%% kapa: exam=1 round=25 year=2014 period=2 form=A subject=accounting subject_ko=회계학 seq=35 number=75
%% source: Q-Net 감정평가사 기출문제 (cst00309) artlSeq=5101207 file=extracted/1008/5101207/2014년 제25회 감정평가사 1차시험 기출문제/2교시(A형).hwp
%% conversion: kordoc 4.24.51 -> tools/gen1.py (auto); review=proofread-sonnet; flags=none
\begin{question}{75}
(주)대한은 단일 종류의 제품을 생산ㆍ판매하고 있다. 20×1년도 단위당 판매가격은 ￦4,000, 단위당 변동원가는 ￦3,500, 연간 총고정원가는 ￦500,000으로 예상된다. 20×1년 중에 특정 고객으로부터 제품 100단위를 구입하겠다는 주문(이하, 특별주문)을 받았다. 특별주문을 수락할 경우 단위당 변동원가 중 ￦500을 절감할 수 있으며, 배송비용은 총 ￦10,000이 추가로 발생한다. 특별주문을 수락하더라도 여유설비가 충분하기 때문에 정상적인 영업활동이 가능하다. (주)대한이 특별주문을 수락하여 ￦30,000의 이익을 얻고자 한다면, 단위당 판매가격을 얼마로 책정해야 하는가?
\begin{choices}
  \item ￦3,100
  \item ￦3,300
  \item ￦3,400
  \item ￦3,500
  \item ￦3,600
\end{choices}
\end{question}
\answer{3}
